\chapter{Theoretical Framework}

The objective of this chapter is to present the theoretical foundation within the field of International Political Economy (IPE) that will serve, in the third chapter, as the analytical basis for interpreting the diffusion of power in the context of the anonymous network The Onion Router (TOR), as reported in journalistic documents.

In this first chapter, we begin from a broad perspective of the traditional concept of power in International Relations (IR) and move toward the alternative conception developed by Susan Strange (1988), which she termed \textit{structural power}. This alternative concept is presented with particular attention to the structure of knowledge. Such structuring aims to provide sufficient theoretical foundations for explaining the phenomenon of \textit{diffusion of power}---introduced subsequently as part of this dissertation's theoretical framework. The phenomenon of \textit{diffusion of power} is grounded in the notion of structural power.

Next, we explore the three dimensions that, according to our interpretation, are essential for understanding the phenomenon of diffusion of power: \textit{authority}, \textit{control}, and \textit{outcomes}. These three dimensions are present in Strange's works (1988, 1996), although they were not explicitly formulated as such by the author. This section will therefore provide the necessary elements for assessing the existence and scope of power diffusion within the TOR network (which is part of the knowledge structure)---an analysis developed in the third chapter of this dissertation.

Finally, we briefly present what the IR literature on power has argued in the cyber domain, focusing on the concepts of \textit{cyber power} and \textit{power diffusion} developed by Joseph Nye (2011).

In summary, this first chapter is organized around the following thematic topics: a brief biographical introduction to Strange; a mapping of traditional debates on power in IR, aimed at highlighting the theoretical gaps explored by Strange---which led to her break with IR literature and the development of her structural analytical framework in IPE; a presentation of structural power and the structure of knowledge; a discussion of the phenomenon of power diffusion; an exploration of the three dimensions---\textit{authority}, \textit{control}, and \textit{outcomes}---and their respective connections with the phenomenon of \textit{power diffusion}; and, finally, an examination of the current IR literature on power in the cyber domain.
